\documentclass[11pt]{article}
\usepackage[margin=2.5cm]{geometry}
\usepackage{amsmath,amssymb,amsthm}
\usepackage[hidelinks]{hyperref}
\newtheorem{theorem}{Theorem}
\newtheorem{lemma}[theorem]{Lemma}
\newtheorem{corollary}[theorem]{Corollary}
\title{No torus with ten convex quadrilateral faces}
\author{Lee Marshall}
\date{29 October 2026}
\begin{document}
\maketitle

\begin{abstract}
% LEE: abstract (Lee writes).
\end{abstract}

\section{Introduction}
% LEE: introduction (Lee writes). Context: Grünbaum--Szilassi 2009, Conjecture 3.2; Seamons 2026 (F=11).

\section{Statement}
We use the definitions of Gr\"unbaum and Szilassi \cite[p.~21]{GS}.
\begin{itemize}
\item Faces are planar simple polygons with no collinear adjacent edges.
\item The faces at a vertex form one circuit.
\item Adjacent faces are not coplanar.
\item The polyhedron is acoptic: every component of the intersection of two faces is a common vertex or a common
  edge.
\end{itemize}

\begin{theorem}\label{thm:main}
No acoptic polyhedral torus in $\mathbb{R}^3$ has exactly ten faces that are all strictly convex quadrilaterals.
\end{theorem}

\begin{theorem}\label{thm:strong}
Let $Q_i=(i,i+1,i+8,i+7)$, indices mod $10$, be the faces of the map $\{4,4\}_{(3,1)}$. The map has no
realisation in $\mathbb{R}^3$ with planar, strictly convex faces and non-coplanar adjacent faces, even if faces
may intersect.
\end{theorem}

\begin{corollary}
Assume the eleven-face witness of \cite{Seamons} is correct, and assume \cite[Prop.~3.1]{GS}. Then an acoptic
torus with $F$ strictly convex quadrilateral faces exists if and only if $F=9$ or $F\ge 11$.
\end{corollary}

\section{Reduction to one map}
\begin{lemma}\label{lem:star}
In a torus as in Theorem~\ref{thm:main}, $V=10$, $E=20$ and every vertex has degree $4$.
\end{lemma}
\begin{proof}[Proof outline]
\begin{itemize}
\item From $4F=2E$ and $V-E+F=0$ we get $E=20$ and $V=10$.
\item Two convex faces have a convex intersection. Being acoptic, it is empty, a common vertex or a common edge.
\item At a vertex $v$ of degree $d$, the points $v$, its $d$ neighbours and its $d$ opposite corners are pairwise
  distinct. A coincidence would make some segment $vw$ a diagonal of one face lying in a second face, which the
  previous point forbids.
\item Hence $2d+1\le 10$, so $d\le 4$. The degree sum is $40=4V$, so every vertex has degree $4$.
\end{itemize}
\end{proof}

\begin{lemma}\label{lem:map}
The face lattice is that of $\{4,4\}_{(3,1)}$.
\end{lemma}
\begin{proof}[Proof outline]
\begin{itemize}
\item A degree-$4$ quadrangulation of the torus is $\mathbb{Z}^2/L$ with $[\mathbb{Z}^2:L]=10$.
\item By Lemma~\ref{lem:star}, $L\cap\{-2,\dots,2\}^2=\{0\}$.
\item In Hermite normal form $L=\langle(a,0),(b,d)\rangle$ with $ad=10$, this leaves only
  $L=\langle(10,0),(3,1)\rangle$ and its mirror image.
\item Two independent exhaustive enumerations confirm this.
\end{itemize}
\end{proof}

\section{Sign conditions and the SAT certificate}
Write $\chi(p_1,\dots,p_4)=\operatorname{sign}\det[1\;p_i]$. Every realisation satisfies the following.
\begin{description}
\item[(GP)] For all $a,b$ and all distinct $y_1,\dots,y_4$, the signs of
  $[aby_1y_2][aby_3y_4]$, $-[aby_1y_3][aby_2y_4]$ and $[aby_1y_4][aby_2y_3]$ are all zero or contain both signs
  (the three-term Pl\"ucker identity).
\item[(PL)] $\chi(Q_i)=0$ for every face.
\item[(AD)] If faces $F$ and $G$ share an edge $ab$, $c$ is another vertex of $F$, and $x,y$ are the vertices of
  $G$ off $ab$, then $\chi(a,b,c,x)=\chi(a,b,c,y)\neq 0$. Reason: $G$ is strictly convex, and its plane meets
  $\operatorname{plane}(F)$ exactly in the line $ab$.
\end{description}

\begin{lemma}\label{lem:sat}
GP, PL and AD for $\{4,4\}_{(3,1)}$ are jointly unsatisfiable.
\end{lemma}
\begin{itemize}
\item The CNF has 420 sign variables and 133\,310 clauses.
\item CaDiCaL proves it unsatisfiable. The LRAT proof (92\,KB) is accepted by the CakeML-verified checker
  \texttt{cake\_lpr} and by a second checker.
\item A separately written encoder (no auxiliary variables) gives a CNF that is also unsatisfiable, with a
  verified LRAT proof.
\item Both encoders are tested on genuine nine- and eleven-face tori: every clause holds for their exact
  chirotopes.
\item Already the 8 faces left after deleting $Q_0$ and $Q_3$ give an unsatisfiable CNF, with verified
  LRAT proofs from both encoders. Inside it, a solver-reported group-minimal core uses 16 AD conditions
  and 37 GP relations; its minimality is not certificate-checked.
\end{itemize}
Lemmas~\ref{lem:star}--\ref{lem:sat} prove Theorem~\ref{thm:main}. Lemma~\ref{lem:sat} does not use
non-intersection, which gives Theorem~\ref{thm:strong}.

\section{Data and verification}
Code, certificates and a one-command checker: \url{% LEE: repository / DOI link
}.
\begin{itemize}
\item The SHA-256 hash of the core CNF is
  \texttt{9afbf70b30fd88a8\allowbreak c95839f163f25f1e\allowbreak e21c89f380faf986\allowbreak a1284b239938f4a2}.
\item The check needs no human-readable proof of Lemma~\ref{lem:sat}. Such a proof is not available.
\item Residual literature risk: Simutis's 1977 thesis was not obtained.
\end{itemize}

\section{Discussion}
% LEE: discussion (Lee writes). Open follow-up: acoptic realisations of {4,4}_(3,1) with non-convex faces.

\section*{AI use}
% LEE: AI-use statement (Lee writes; same style as the README disclosure).

\begin{thebibliography}{9}
\bibitem{GS} B.~Gr\"unbaum and L.~Szilassi, Geometric realizations of special toroidal complexes,
  \emph{Contrib.\ Discrete Math.} 4(1) (2009) 21--39, doi:10.55016/ojs/cdm.v4i1.61986.
\bibitem{G99} B.~Gr\"unbaum, Acoptic polyhedra, in \emph{Advances in Discrete and Computational Geometry},
  Contemp.\ Math.\ 223 (1999) 163--199.
\bibitem{Seamons} J.~W.~G.~Seamons, A torus with eleven convex quadrilateral faces, Zenodo (2026),
  doi:10.5281/zenodo.23750303.
\bibitem{BK} U.~Brehm and W.~K\"uhnel, Equivelar maps on the torus, \emph{European J.\ Combin.} 29 (2008)
  1843--1861.
\bibitem{cake} Y.~K.~Tan, M.~J.~H.~Heule and M.~O.~Myreen, Verified propagation redundancy and compositional UNSAT
  checking in CakeML, \emph{Int.\ J.\ Softw.\ Tools Technol.\ Transf.} 25 (2023) 167--184.
\end{thebibliography}
\end{document}
